\chapter{Introducción}
\label{Introducción}

La computación paralela y distribuida es una rama de la informática que se encarga de estudiar modelos y sistemas que permiten ejecutar tareas de forma simultánea o repartidas, entre distintas máquinas con múltiples unidades de procesamiento.
En específico, la computación paralela busca obtener máximo rendimiento para dichas tareas dividiendo los problemas en partes que se puedan ejecutar al mismo tiempo en distintas máquinas.
En cambio, la computación distribuida se enfoca en los en los sistemas cuyos componentes se encuentran conectados en computadoras conectadas en red, comunicándose y coordinándose mediante el paso de mensajes.

En este informe se tratará el tema, y específicamente sobre la aplicación Steam, app. creada por Valve y que es una de las más populares y grandes del mundo del gaming. Se verán las conectividades entre sus distintos servidores, como los servidores encargados del servicio de mensajería de chat, los servidores encargados de las bibliotecas de videojuegos de los usuarios, entre otros. Esto se verá reflejado en distintos modelos, tales como modelos físicos, modelos arquitectónicos y modelos fundamentales.